\documentclass[10pt,reqno]{amsart}

\usepackage[T1]{fontenc}
\usepackage{amsmath}
\usepackage{amssymb}
\usepackage{mathtools}
\usepackage{booktabs}
\usepackage{array}
\usepackage{enumitem}
\usepackage{microtype}
\usepackage{hyperref}
\usepackage{url}
\usepackage{xcolor}

\hypersetup{
  hidelinks,
  pdftitle={Topological Semantics for Scoped Computational Paths},
  pdfauthor={Arthur Freitas Ramos, Ruy J. G. B. de Queiroz, Anjolina Grisi de Oliveira, and Tiago M. L. de Veras}
}

\emergencystretch=3em
\allowdisplaybreaks

\theoremstyle{definition}
\newtheorem{definition}{Definition}[section]
\newtheorem{example}[definition]{Example}
\newtheorem{construction}[definition]{Construction}
\theoremstyle{plain}
\newtheorem{theorem}[definition]{Theorem}
\newtheorem{lemma}[definition]{Lemma}
\newtheorem{proposition}[definition]{Proposition}
\newtheorem{corollary}[definition]{Corollary}
\theoremstyle{remark}
\newtheorem{remark}[definition]{Remark}

\newcommand{\Path}{\mathsf{Path}}
\newcommand{\Trace}{\mathsf{Trace}}
\newcommand{\Step}{\mathsf{Step}}
\newcommand{\RwEq}{\mathrel{\simeq_{\!\mathcal P}}}
\newcommand{\PiOne}{\Pi_1}
\newcommand{\Fin}{\mathrm{fin}}
\newcommand{\Ord}{\mathrm{ord}}
\newcommand{\Top}{\mathsf{Top}}
\newcommand{\Z}{\mathbb Z}
\newcommand{\N}{\mathbb N}
\newcommand{\Circle}{\mathbb R/\mathbb Z}
\urlstyle{tt}
\newcommand{\LeanName}[1]{\path{#1}}

\title[Topological Semantics for Scoped Computational Paths]
  {Topological Semantics for Scoped Computational Paths}

\author{Arthur Freitas Ramos}
\address{Microsoft, USA}
\email{arfreita@microsoft.com}

\author{Ruy J. G. B. de Queiroz}
\address{Centro de Inform\'atica, Universidade Federal de Pernambuco, Brazil}
\email{ruy@cin.ufpe.br}

\author{Anjolina Grisi de Oliveira}
\address{Centro de Inform\'atica, Universidade Federal de Pernambuco, Brazil}
\email{ago@cin.ufpe.br}

\author{Tiago M. L. de Veras}
\address{Departamento de Matem\'atica, Universidade Federal Rural de Pernambuco, Brazil}
\email{tiago.veras@ufrpe.br}

\renewcommand{\shortauthors}{Ramos et al.}

\subjclass[2020]{Primary 03B38, 18N10, 54B30; Secondary 55Q05, 68V15}
\keywords{computational paths, rewriting, fundamental groupoid, quotient topology, topological groupoid, formalized mathematics}

\begin{document}

\raggedbottom

\begin{abstract}
We give a topological semantics for a genuinely scoped computational-path
rewrite presentation.  A presentation consists of endpoint-changing geometric
steps, coherent path realizations, and named endpoint-fixed rewrite rules.  The
rewrite relation is the least equivalence congruence generated by those rules
and the explicit groupoid coherences; it is not ambient equality normalization.

For every presentation we construct a quotient arrow space and a canonical
final-domain topological groupoid.  Its multiplication is continuous on the
quotient of the explicit composable-trace carrier, without any product-quotient
assumption.  We then prove an exact four-way classification of when this final
domain agrees with the ordinary pullback topology.  Thus the main groupoid
theorem is unconditional, while the ordinary-pullback upgrade is a theorem
about a separately checkable topological property, not an assumption hidden in
the construction.

We also prove a positive ordinary-pullback theorem: compactness of the final
composable domain and Hausdorffness of the ordinary domain force the two
topologies to agree.  A discrete finite-presentation corollary supplies an
explicit class of presentations to which the ordinary upgrade applies.

The realization map to endpoint-fixed geometric homotopy classes is a
continuous groupoid morphism.  It is faithful exactly when the presentation is
geometrically complete, and geometric completeness upgrades the comparison to
a homeomorphism on arrow spaces.  The universal presentation, whose rules are
all endpoint-fixed homotopies, recovers the quotient-topologized fundamental
groupoid.  A concrete presentation on the unit additive circle has explicit
zero, addition, and reversal rules, an inductive integer normal form, a
continuous injection of the winding loop quotient, and a proved equivalence of
that loop carrier with $\Z$.  The actual torus $(\mathbb R/\mathbb Z)^2$ is
then classified by a two-coordinate winding theorem.  The Lean development is
an appendix and artifact: the mathematics below is independent of Lean.
\end{abstract}

\maketitle

\section{Introduction}

Computational paths were introduced to make equality witnesses explicit as
traces of computational steps.  The authors' program has developed this idea
from propositional equality and identity types through path normalization,
groupoid structure, and topological applications
\cite{RamosDeQueirozDeOliveira2016propositional,Ramos2017identity,
RamosDeQueirozDeOliveiraDeVeras2018explicit,
DeVerasRamosDeQueirozDeOliveira2025topological,
DeVerasRamosDeQueirozDeOliveira2025weak}.  This viewpoint is useful only while the trace
syntax remains visible: if every proof of the same ambient equality is silently
identified, then the intended computational information has already been
discarded.  The central design problem is therefore to define a rewrite
system whose semantic interpretation is sound without making semantic equality
the rewrite relation itself.

The broader publication record includes the thesis and thesis abstract,
the transport and univalence presentation, the 2021 calculations of
fundamental groups and groupoids, and the Portuguese-language workshop paper
on term conversion and homotopy \cite{Ramos2018explicitThesis,Ramos2019thesisAbstract,
Ramos2017transport,RamosDeQueirozDeOliveira2021calculation,
RamosDeQueirozDeOliveira2021fundamental,
RamosDeQueirozDeOliveira2021conversao}.  The authors' book gives the broader
calculus and its historical context \cite{RamosDeQueirozDeOliveiraGabbay2026}.
The earlier EBL and WoLLIC papers, the Springer monograph, and the
fundamental-groups thesis are also part of that record
\cite{RamosDeQueirozDeOliveira2014,RamosDeQueirozDeOliveira2015,
RamosDeQueirozDeOliveira2016,RamosDeQueirozDeOliveira2017,
deQueirozDeOliveiraRamos2016,Ramos2019thesis}.

This paper develops the topological version of that design.  We work with
proof-irrelevant ambient paths, meaning ordinary continuous interval paths
together with endpoint equalities, and with explicit finite geometric traces.
The word ``computational'' refers to the trace and its declared rewrites; it
does not refer to homotopy type theory, univalence, or a higher-inductive
identity type.  In particular, our fundamental-groupoid comparison is a
semantic map out of a presented trace quotient, not an identification of
proof-relevant identity types with topological spaces \cite{HofmannStreicher98,
RamosDeQueirozDeOliveira2016,RamosDeQueirozDeOliveira2017,
RamosDeQueirozDeOliveira2015}.

The contributions are the following.

\begin{enumerate}[label=(\roman*),leftmargin=2.2em]
\item We define continuous geometric rewrite presentations and a scoped rewrite
  relation generated by named rules, groupoid operations, and congruence.
  Soundness for geometric realization is proved by induction over the complete
  derivation system.
\item We construct the quotient arrow space, its endpoint maps, identities,
  reversal, and a canonical final-domain topological groupoid.  The final
  composable domain is the quotient of explicit composable traces, so the
  continuity proof does not use an unproved product of quotient maps.
\item We prove an exact classification of the ordinary-pullback upgrade.  Four
  conditions are equivalent: quotientness of the ordinary composable-pair map,
  continuity of the inverse final-to-ordinary identification, equality of the
  two composable-pair topologies, and quotientness of the raw ordinary-pair
  map.  An open-map criterion is proved as a sufficient theorem.
\item We prove a positive compact-Hausdorff compatibility theorem and its
  discrete finite-presentation corollary.  Thus the paper contains an
  instantiated ordinary-pullback class, not only an exact obstruction.
\item We prove an effective normal-form criterion: a scoped normalizer that is
  geometrically separating yields geometric completeness.  The criterion is
  instantiated by the finite-generator circle calculation.
\item We prove functoriality for continuous presentation maps, a realization
  comparison and completeness criterion, and a universal presentation theorem.
  The comparison preserves the final composable-domain operation, not only a
  set-level map on arrows.
\item We work out the actual additive circle from one oriented generator, with
  cancellation rules and a derived integer normal form.  We then compute the
  genuine topological torus loop quotient as $\Z\times\Z$ by the product of
  the two circle covering classifications.  Neither calculation makes a
  discrete-topology claim about the geometric loop spaces.
\end{enumerate}

The paper is deliberately separate from the authors' companion manuscript,
``The Seifert--van Kampen Theorem via Computational Paths,'' which is currently
an unpublished arXiv preprint rather than a published paper
\cite{RamosEtAl2025SVKPreprint}.  That preprint studies a specific
fundamental-group calculation.  Here the principal result is the general
topological construction for scoped rewrite presentations, together with the
ordinary/final topology boundary and the universal semantic completion.  The
fundamental-groupoid terminology follows the standard topological treatment
of groupoids and covering-space calculations \cite{Brown06,Hatcher02}.

To make the novelty boundary explicit, the earlier computational-path papers
are used as foundations and comparisons, not silently counted as results of
this manuscript.  The 2016--2018 work develops the identity-type and explicit
trace calculus; the 2021 papers calculate selected fundamental groups and
groupoids; the accepted topological application and the weak-groupoid paper
provide adjacent geometric and higher-structural results.  The unpublished
SVK manuscript is a companion preprint, not a published source and not a
result claimed here.  The present paper contributes the scoped topological
quotient, the final/ordinary domain theorem, the positive compatibility class,
the effective completeness criterion, and the circle--torus applications.

\begin{table}[ht]
\centering
\small
\setlength{\tabcolsep}{3pt}
\begin{tabular}{>{\raggedright\arraybackslash}p{0.29\linewidth}>{\raggedright\arraybackslash}p{0.30\linewidth}>{\raggedright\arraybackslash}p{0.33\linewidth}}
\toprule
Record & Role in this paper & Result not imported from the record\\
\midrule
Identity type and explicit paths
  \cite{RamosDeQueirozDeOliveira2016propositional,Ramos2017identity,
  RamosDeQueirozDeOliveiraDeVeras2018explicit}
  & Syntax, equality background, and groupoid motivation
  & Scoped topological quotient and its two composable-domain topologies\\
2021 computational-path calculations
  \cite{RamosDeQueirozDeOliveira2021calculation,
  RamosDeQueirozDeOliveira2021fundamental,
  RamosDeQueirozDeOliveira2021conversao}
  & Earlier concrete algebraic and homotopical applications
  & Compact-Hausdorff compatibility and normal-form completeness\\
Topological and weak-groupoid papers
  \cite{DeVerasRamosDeQueirozDeOliveira2025topological,
  DeVerasRamosDeQueirozDeOliveira2025weak}
  & Adjacent topological and higher-categorical context
  & Final-domain semantics and the exact ordinary/final comparison\\
SVK companion preprint
  \cite{RamosEtAl2025SVKPreprint}
  & Separate unpublished calculation
  & No theorem here depends on its proof or is presented as its result\\
This paper
  & Standalone topological semantics
  & Scoped quotient, positive topology class, effective completeness, circle,
    and genuine torus classification\\
\bottomrule
\end{tabular}
\caption{Theorem-level novelty and overlap boundary.}
\label{tab:novelty}
\end{table}

\section{Geometric traces and scoped presentations}

Throughout, $X$ is a topological space.  Let $E$ be a topological space of
primitive steps, with continuous endpoint maps $s,t:E\to X$ and a continuous
family of realized interval paths
\[
  \rho(e):s(e)\leadsto t(e).
\]
The continuity requirement means that the corresponding map into the
compact-open path space $X^I$ is continuous.  The endpoint-indexed notation
below is convenient for stating composition, but the underlying trace syntax
is independent of a proof of ambient equality.

\begin{definition}[geometric trace]
The geometric trace inductive family is generated by
\[
\frac{}{\Trace(a,a)},\qquad
\frac{e:E}{\Trace(s(e),t(e))},\qquad
\frac{p:\Trace(a,b)\quad q:\Trace(b,c)}
  {p\mathbin{; }q:\Trace(a,c)}.
\]
\[
\frac{p:\Trace(a,b)}{p^{-1}:\Trace(b,a)}.
\]
Its realization is defined recursively by constant paths, $\rho(e)$,
concatenation, and reversal.  The trace length $\ell(p)\in\N$ is defined by
length $0$ for reflexivity, length $1$ for a primitive step, addition under
concatenation, and preservation under reversal.
\end{definition}

The topology on each trace fiber observes both $\ell(p)$ and the realized
path.  A coherent open path is a triple
\[
  u=(p,\gamma,H),
\]
where $p:\Trace(a,b)$, $\gamma:a\leadsto b$ is an arbitrary continuous
interval path, and $H$ is an endpoint-fixed homotopy from $\gamma$ to the
realization of $p$.  Let $T$ denote the endpoint-varying carrier of coherent
open paths.  Its topology observes source, target, trace length, realized
trace, and chosen geometric representative.  This choice is important: the
topology remembers geometric motion while the quotient relation below controls
which computational traces are identified.

\begin{definition}[scoped presentation]
A continuous geometric rewrite presentation $\mathcal P$ on $(X,E)$ consists
of a predicate
\[
  R_{\mathcal P}(p,q),\qquad p,q:\Trace(a,b),
\]
and, for every instance of $R_{\mathcal P}(p,q)$, an endpoint-fixed homotopy
between the realizations of $p$ and $q$.
\end{definition}

The semantic witness is part of the presentation data.  No completeness
condition is included in the definition.

\begin{definition}[scoped rewrite equality]
For parallel traces, write $p\RwEq q$ for the inductively generated relation
whose constructors are reflexivity, the named generators $R_{\mathcal P}$,
symmetry, transitivity, congruence under concatenation and reversal, and the
following strict groupoid coherences:
\[
\begin{array}{lll}
(1_a);p\RwEq p, & p;(1_b)\RwEq p, & (p;q);r\RwEq p;(q;r),\\
p^{-1};p\RwEq 1_b, & p;p^{-1}\RwEq 1_a, & (p^{-1})^{-1}\RwEq p,\\
(1_a)^{-1}\RwEq 1_a, & (p;q)^{-1}\RwEq q^{-1};p^{-1}. &
\end{array}
\]
Only the rules declared by $\mathcal P$ and these structural constructors are
available.  In particular, ambient equality of traces is not a constructor.
\end{definition}

\begin{theorem}[realization soundness]
If $p\RwEq q$, then the realized paths are endpoint-fixed homotopic.
\end{theorem}

\begin{proof}
Induct on the derivation of $p\RwEq q$.  Reflexivity is the constant
homotopy, and a named generator is sound by the corresponding presentation
witness.  Symmetry and transitivity use symmetry and vertical composition of
endpoint-fixed homotopies.  For concatenation, the two inductive homotopies
are combined by horizontal concatenation; for reversal, reverse the homotopy
parameter.  The unit, associativity, cancellation, double-reversal, and
reversal-of-composite cases are the standard endpoint-fixed homotopies of
interval paths.  Each case is therefore discharged by the corresponding
constructor of the ambient homotopy equivalence relation.
\end{proof}

\begin{remark}
The soundness theorem is one-way.  Equality of realized geometric homotopy
classes does not imply scoped rewrite equality unless a separate geometric
completeness theorem is proved.  This distinction is the logical reason that
the presentation can retain computational information.
\end{remark}

\section{The quotient arrow space}

The scoped relation on endpoint-varying coherent paths compares endpoints and
then compares their traces using the declared presentation.  Write $u\approx
v$ for this relation on $T$.  The quotient arrow space is
\[
  G_{\mathcal P}=T/{\approx},
\]
with the quotient topology induced by $q:T\to G_{\mathcal P}$.  Endpoint
maps descend to continuous maps
\[
  s,t:G_{\mathcal P}\longrightarrow X.
\]
Reflexive paths and reversal descend because the scoped relation is closed
under the corresponding constructors.

The explicit composable carrier is
\[
  T^{(2)}=\{(u,v)\in T\times T:t(u)=s(v)\}.
\]
Its quotient relation is componentwise scoped equivalence, and its quotient is
denoted $G_{\mathcal P}^{(2),\Fin}$.  Concatenation of representatives gives a
well-defined map
\[
  m_{\Fin}:G_{\mathcal P}^{(2),\Fin}\longrightarrow G_{\mathcal P}.
\]
The quotient topology is chosen precisely so that this map has the required
continuity.

\begin{theorem}[canonical final-domain groupoid]
For every continuous geometric rewrite presentation $\mathcal P$:
\begin{enumerate}[label=(\alph*),leftmargin=2.2em]
\item $s,t$, the identity map $X\to G_{\mathcal P}$, and reversal
  $G_{\mathcal P}\to G_{\mathcal P}$ are continuous;
\item $m_{\Fin}$ is continuous;
\item the descended operations satisfy the left and right unit laws, inverse
  laws, and associativity.
\end{enumerate}
Thus $G_{\mathcal P}$ is a final-domain topological groupoid.  The phrase
``final-domain topological groupoid'' records the domain on which
multiplication is continuous.
\end{theorem}

\begin{proof}
Continuity of the endpoint maps and identities follows from the quotient
criterion, because the corresponding raw maps are continuous.  The same
argument applies to reversal.  The raw concatenation map
$T^{(2)}\to T$ is continuous and respects the componentwise quotient
relation, so it descends to $m_{\Fin}$; continuity follows from the defining
quotient topology on $G_{\mathcal P}^{(2),\Fin}$.

For the algebraic laws, choose representatives.  Left and right units reduce
to the scoped constructors $(1_a);p\RwEq p$ and $p;(1_b)\RwEq p$.  The two
inverse laws reduce to $p^{-1};p\RwEq 1_b$ and $p;p^{-1}\RwEq 1_a$.
Associativity reduces to $(p;q);r\RwEq p;(q;r)$.  The quotient soundness
criterion then turns each scoped derivation into equality of quotient arrows.
The result is independent of representative choice because concatenation was
defined through the componentwise quotient.
\end{proof}

\section{Ordinary and final composable-pair topologies}

There is another natural composable domain.  Let
\[
  G_{\mathcal P}^{(2),\Ord}
  =\{(x,y)\in G_{\mathcal P}\times G_{\mathcal P}:t(x)=s(y)\}
\]
with the subspace topology.  The underlying set of this space is the same as
the underlying set of the final domain, but its topology need not be the same.
The canonical map
\[
  q^{(2)}:T^{(2)}\longrightarrow G_{\mathcal P}^{(2),\Ord}
\]
is continuous and surjective.  Product quotient maps are not quotient maps in
arbitrary topological spaces, so quotientness of $q^{(2)}$ is an additional
property, not a formal consequence of quotientness of $q$.

\begin{definition}[product-quotient compatibility]
The presentation has product-quotient compatibility if $q^{(2)}$ is a quotient
map onto the ordinary composable-pair space.
\end{definition}

The final domain can be identified with the ordinary domain by the bijection
that is the identity on the underlying set.  One direction is always
continuous; the reverse direction is the exact point at which compatibility is
needed.

\begin{theorem}[ordinary/final comparison]
For every presentation $\mathcal P$, the following are equivalent:
\begin{enumerate}[label=(\arabic*),leftmargin=2.2em]
\item product-quotient compatibility holds;
\item the identity map from the ordinary composable domain to the final
  composable domain is continuous;
\item the ordinary and final composable-pair topologies coincide;
\item the raw ordinary-pair map, viewed from $T^{(2)}$, is a quotient map.
\end{enumerate}
When these conditions hold, multiplication is continuous from the ordinary
pullback domain.  No condition in this theorem is used by the canonical
final-domain groupoid theorem.
\end{theorem}

\begin{proof}
The final topology is the coinduced topology for $q^{(2)}$; hence quotientness
of $q^{(2)}$ is equivalent to equality between the final topology and the
ordinary topology.  Equality of the two topologies is equivalent to continuity
of the identity in the reverse direction because the forward identity is
always continuous.  The raw ordinary-pair map is the composite of the raw
composable quotient with the identity between the two pair domains.  The
composition and cancellation criteria for quotient maps therefore show that
its quotientness is equivalent to quotientness of $q^{(2)}$.  Finally, when
$q^{(2)}$ is quotient, multiplication is continuous by the quotient
criterion applied to its continuous representative-level concatenation.
\end{proof}

\begin{proposition}[open-map criterion]
If the bijection from the final composable quotient to the ordinary pair space
is an open map, then product-quotient compatibility holds.  In particular, an
open quotient presentation of the composable domain supplies the ordinary
pullback upgrade.
\end{proposition}

\begin{theorem}[compact-Hausdorff compatibility]
Suppose that the final composable domain
$G_{\mathcal P}^{(2),\Fin}$ is compact and that the ordinary composable domain
$G_{\mathcal P}^{(2),\Ord}$ is Hausdorff.  Then product-quotient compatibility
holds.  Hence the final and ordinary composable-pair topologies agree and
ordinary-pullback multiplication is continuous.
\end{theorem}

\begin{proof}
The identity map

\[
  j:G_{\mathcal P}^{(2),\Fin}\longrightarrow
  G_{\mathcal P}^{(2),\Ord}
\]

is a continuous bijection by the comparison theorem.  A continuous bijection
from a compact space to a Hausdorff space is a homeomorphism: its image of
every closed set is compact, hence closed, so its inverse is continuous.  The
quotient map defining the final domain followed by $j$ is therefore the raw
ordinary-pair map followed by a homeomorphism, and is a quotient map.  This is
exactly product-quotient compatibility.
\end{proof}

\begin{corollary}[finite discrete presentations]
Let $\mathcal P$ be a scoped presentation whose arrow space is finite and
discrete and whose explicit composable-trace carrier is discrete.  Then the
ordinary and final composable-pair topologies coincide.
\end{corollary}

\begin{proof}
The quotient of a discrete space is discrete for the quotient topology, so
the final composable domain is discrete.  It is finite because its underlying
set is a subset of the square of the finite arrow space.  The ordinary
composable domain is a finite subspace of a discrete square and is therefore
Hausdorff.  The compact-Hausdorff theorem applies.
\end{proof}

The compact-Hausdorff theorem is the promised positive class: it establishes
the ordinary-pullback upgrade from independently checkable properties of the
actual presentation, rather than leaving the four-way criterion as the only
result.  The final-domain theorem remains the general statement and does not
need these hypotheses.

\section{Functoriality}

Let $\mathcal P$ and $\mathcal Q$ be presentations on $(X,E)$ and $(Y,F)$.
A presentation map consists of a continuous map $f:X\to Y$, a continuous
step map preserving endpoints and realizations, and a proof that every named
rule of $\mathcal P$ maps to a scoped derivation in $\mathcal Q$.

\begin{theorem}[functoriality]
Every presentation map $F:\mathcal P\to\mathcal Q$ induces a continuous map
\[
  F_1:G_{\mathcal P}\longrightarrow G_{\mathcal Q}
\]
that preserves source, target, identities, reversal, and final-domain
composition.  The assignment preserves identity maps and composites.
\end{theorem}

\begin{proof}
Induction on scoped rewrite derivations proves that the step map sends related
traces to related traces.  It therefore descends through the quotient.  The
raw map on coherent paths is continuous, so the quotient criterion gives
continuity on arrows.  Applying the step map to reflexivity, reversal, and
concatenation gives the corresponding preservation equations before quotienting.
The same argument on the explicit composable carrier gives a continuous map on
final composable domains and proves final composition preservation.  Identity
and composite step maps are definitionally the identity and composite maps on
raw traces; quotient extensionality completes the proof.
\end{proof}

The same construction gives ordinary-pair continuity for every presentation
map.  If both source and target presentations satisfy product-quotient
compatibility, the ordinary multiplication maps are continuous as well.  This
is a derived theorem, not the definition of functoriality.

\section{Geometric comparison and completeness}

Let $H_{\mathrm{geo}}$ be the quotient of coherent open paths by equality of
endpoints and endpoint-fixed homotopy of chosen geometric representatives.
The realization comparison
\[
  R_{\mathcal P}:G_{\mathcal P}\longrightarrow H_{\mathrm{geo}}
\]
maps a scoped class to its geometric homotopy class.

\begin{theorem}[comparison morphism]
The map $R_{\mathcal P}$ is continuous, surjective on the geometric quotient,
and preserves source, target, identity, reversal, and final-domain composition.
In particular it is a continuous groupoid morphism for the final-domain
structures.
\end{theorem}

\begin{proof}
Realization soundness makes the map well-defined.  Its composite with the raw
quotient map is the continuous geometric quotient map, so the quotient
criterion proves continuity.  Every geometric class has a coherent raw
representative by definition, proving surjectivity.  The operation equations
follow from the recursive realization of reflexivity, concatenation, and
reversal.  For the final domain, map a strong pair to the corresponding strong
pair of geometric classes; the quotient representative of its composition is
the geometric composition of the two component representatives.
\end{proof}

\begin{definition}[geometric completeness]
The presentation is geometrically complete if, for all raw coherent paths
$u,v$, equality of their geometric quotient codes implies scoped equivalence:
\[
  R_{\mathcal P}(q(u))=R_{\mathcal P}(q(v))
  \quad\Longrightarrow\quad u\approx v.
\]
\end{definition}

\begin{theorem}[faithfulness and homeomorphism]
The comparison $R_{\mathcal P}$ is faithful if and only if $\mathcal P$ is
geometrically complete.  If $\mathcal P$ is geometrically complete, then
\[
  R_{\mathcal P}:G_{\mathcal P}\xrightarrow{\ \cong\ }H_{\mathrm{geo}}
\]
is a homeomorphism on arrow spaces and an isomorphism of abstract groupoids.
\end{theorem}

\begin{proof}
If the comparison is injective, equality of geometric codes for two raw paths
gives equality of their scoped quotient classes, which is precisely scoped
completeness.  Conversely, completeness turns equality of comparison images
into equality of scoped classes.  The comparison is always surjective and
continuous.  Under completeness, define the inverse by lifting a geometric
class to a raw representative and taking its scoped class.  Completeness makes
this lift well-defined, and the quotient criterion makes it continuous.  The
two composites are identity maps by quotient induction.
\end{proof}

\begin{definition}[effective normal-form certificate]
An effective normal-form certificate for $\mathcal P$ is a map

\[
  N:T\longrightarrow T
\]

such that, for all raw coherent paths $u,v$,

\begin{enumerate}[label=(\alph*),leftmargin=2.2em]
\item $u\approx N(u)$;
\item equality of geometric quotient codes implies $N(u)=N(v)$; and
\item $N(u)=N(v)$ implies $N(u)\approx N(v)$.
\end{enumerate}

The first clause is a derivable normalization certificate, the second is
semantic separation of normal forms, and the third records the scoped proof
that equal normal forms represent the same presented arrow.
\end{definition}

\begin{theorem}[normal-form completeness criterion]
An effective normal-form certificate implies geometric completeness.  Therefore
the realization comparison is faithful and is a homeomorphism on arrow spaces.
\end{theorem}

\begin{proof}
Assume that the geometric codes of $u$ and $v$ are equal.  Clause (b) gives
$N(u)=N(v)$, and clause (c) gives $N(u)\approx N(v)$.  Combining this with
clause (a) for $u$ and the symmetric form of clause (a) for $v$ yields

\[
  u\approx N(u)\approx N(v)\approx v.
\]

This is precisely geometric completeness.  The faithfulness and homeomorphism
claims follow from the comparison theorem.
\end{proof}

\section{The universal geometric presentation}

For every topological space $X$, take the primitive step space to be the
continuous interval path space $X^I$.  Its endpoint maps are evaluation at
$0$ and $1$.  Declare two parallel traces to be named generators exactly when
their realized paths are endpoint-fixed homotopic.

\begin{theorem}[universal completion]
For the universal presentation $\mathcal U_X$, scoped rewrite equality of
parallel traces is equivalent to endpoint-fixed homotopy of their realizations.
Consequently,
\[
  G_{\mathcal U_X}\cong \Pi_1^q(X)
\]
as abstract groupoids, and the comparison is a homeomorphism on arrow spaces.
\end{theorem}

\begin{proof}
Soundness is the general soundness theorem.  Conversely, an endpoint-fixed
homotopy is itself one of the named universal generators, so it is a scoped
rewrite in one step.  The equivalence of quotient relations follows after
including endpoint equality.  The geometric completeness theorem then gives
the homeomorphism and groupoid isomorphism.
\end{proof}

\begin{remark}
The multiplication in this completion is understood on the final composable
domain.  This is the unconditional topological statement.  For compact final
composable domains with Hausdorff ordinary domains, the compact-Hausdorff
theorem above identifies this multiplication with the ordinary pullback
operation.
\end{remark}

\section{The additive circle: a finite generator and its completion}

Let $\mathbb T=\mathbb R/\mathbb Z$ be the unit additive circle.  We use the
actual topological circle, not an inductive one-point model.  The minimal
presentation has one oriented primitive step $a$ at the base point $0$, with
realization
\[
  a(t)=[t]\in\mathbb T.
\]
The formal reversal supplies $a^{-1}$; no integer-indexed primitive step is
assumed.  A based trace is therefore a finite signed word in $a$ and
$a^{-1}$.  The only non-structural generators are the two cancellation rules
\[
  a;a^{-1}\RwEq 1_0,
  \qquad a^{-1};a\RwEq 1_0,
\]
closed under the scoped congruence.

For a signed word $w$, let $\operatorname{exp}(w)$ be the number of positive
letters minus the number of negative letters.  Let $c_n$ be the canonical
word consisting of $n$ copies of $a$ when $n\geq 0$ and $-n$ copies of
$a^{-1}$ when $n<0$.

\begin{lemma}[finite-generator reduction]
Every signed word $w$ has a finite scoped derivation
\[
  w\RwEq c_{\operatorname{exp}(w)}.
\]
Moreover, two canonical words are scoped equivalent only when their
exponents are equal after geometric realization.
\end{lemma}

\begin{proof}
Scan the word from left to right while retaining a reduced signed word.  If
the next sign agrees with the last retained sign, append it.  If the signs
differ, delete the adjacent pair; this is one of the two cancellation
generators, possibly whiskered by the already scanned prefix and the
unscanned suffix.  Each deletion strictly decreases the word length, so the
procedure terminates.  A terminal reduced word has no adjacent sign change
and hence consists entirely of positive letters or entirely of negative
letters.  The integer $\operatorname{exp}$ is unchanged by every deletion,
so the terminal word is exactly $c_{\operatorname{exp}(w)}$.  The same
invariant is the endpoint of the lift to $\mathbb R$; therefore endpoint-fixed
homotopy of canonical words forces equality of their exponents.
\end{proof}

The reduction lemma is the effective normal-form certificate from
Section~4: the reduction derivation gives clause (a), the lifted endpoint gives
clause (b), and equality of canonical words gives clause (c).  Thus the
minimal circle presentation is geometrically complete on its based loop
fiber.

For comparison with the checked artifact, the completed integer presentation
has a primitive step $\lambda_n$ for each $n\in\Z$,
\[
  \lambda_n(t)=[tn].
\]
Its named rules are the derived integer rules
\[
  (n;m)\RwEq(n+m),\qquad 1_0\RwEq 0,
  \qquad (n)^{-1}\RwEq(-n).
\]
The finite-generator argument above explains why this completion does not
build the integer classification into the primitive syntax.

Let $L$ be the quotient of coherent based circle loops by endpoint-fixed
homotopy.  The map sending a loop class to its winding number and the map
sending $n$ to the standard loop class are inverse maps.

\begin{theorem}[circle winding classification]
The finite-generator scoped presentation is geometrically complete on the
based loop fiber, and there is an equivalence
\[
  L\simeq\Z.
\]
The induced continuous map from $L$ into the scoped arrow carrier is
injective, and its image is exactly the source-target fiber over the base
point.  In particular the scoped presentation is nondegenerate: the images of
the zero and one winding loops are distinct.
\end{theorem}

\begin{proof}
The quotient map to winding is well-defined because endpoint-fixed homotopies
lift to homotopies with equal lifted endpoints.  The standard loop of winding
$n$ has lift $t\mapsto tn$.  Every based loop is homotopic to the standard loop
with the endpoint of its lift, so the two maps are inverse.  The scoped map is
continuous by the quotient criterion.  If two of its values agree, soundness
gives a geometric homotopy between the two chosen loops; winding invariance
then gives equality of their integer normal forms.  The image condition is
obtained by quotient induction on an arrow and by inserting the based loop
representative when the endpoints are both $0$.  The zero and one classes are
distinct because winding sends them to distinct integers.
\end{proof}

The Lean declarations \LeanName{circleBasedNormalForm},
\LeanName{circleBasedNormalForm_scoped}, and
\LeanName{circleBasedNormalFormCertificate} record
the same based completeness argument for the completed integer presentation.
The equivalence $L\simeq\Z$ is a quotient classification, not a claim that
$L$ or the corresponding scoped fiber is discrete.

\section{The genuine topological torus}

Let
\[
  \mathbb T^2=(\mathbb R/\mathbb Z)\times(\mathbb R/\mathbb Z)
\]
with base point $(0,0)$.  Use two finite oriented generators
\[
  a(t)=([t],[0]),
  \qquad b(t)=([0],[t]).
\]
The scoped rules are the inverse cancellations for $a$ and $b$, together with
the commuting square
\[
  a;b\RwEq b;a.
\]
The commuting square is sound: it is the projection modulo $\mathbb Z^2$ of
the evident square homotopy in $\mathbb R^2$.
Closure under symmetry, reversal, and whiskering supplies the corresponding
commuting rules for all four sign combinations.

For a signed word $w$, let $m(w)$ and $n(w)$ be its signed counts of $a$ and
$b$.  Move every $b$-letter to the right using the commuting square and then
perform the two cancellation algorithms from the circle.  This is a finite
terminating derivation
\[
  w\RwEq a^{m(w)};b^{n(w)}.
\]
The pair $(m(w),n(w))$ is invariant under every scoped rule.

\begin{theorem}[torus winding classification]
For the actual topological torus, the quotient $L_{\mathbb T^2}$ of based
coherent loops by endpoint-fixed homotopy satisfies
\[
  L_{\mathbb T^2}\cong\Z\times\Z.
\]
The finite-generator scoped presentation is geometrically complete on the
based loop fiber, and its realization comparison is injective there.
\end{theorem}

\begin{proof}
The universal cover is the product map
\[
  \mathbb R^2\longrightarrow\mathbb T^2,
  \qquad (x,y)\longmapsto([x],[y]).
\]
The lift of a based loop has endpoint $(m,n)\in\mathbb Z^2$, and endpoint-fixed
homotopy preserves this pair.  Conversely, the two coordinate projections of
any based torus loop are circle loops; the circle lifting theorem gives
endpoint-fixed homotopies to the standard loops of windings $m$ and $n$.
Taking their product gives a homotopy to $a^m;b^n$.  Thus the coordinate
winding pair classifies geometric loops.  The word reduction above gives the
corresponding scoped derivation, so the normal-form criterion proves geometric
completeness and injectivity.
\end{proof}

This example is genuinely topological: the two coordinates are realized by
the additive-circle covering map, and the product homotopy is proved at the
level of continuous interval paths.  It is not the older synthetic torus
carrier and does not use the companion SVK preprint.  The Lean artifact records
the product-winding certificate in
\LeanName{TopologicalTorus.certificate}; as in the
circle case, the final-domain topological groupoid theorem is unconditional,
while no discrete-topology claim is made for the geometric torus loop space.

\section{Relation to logic and rewriting}

The logical content of the construction lies in the boundary between syntax
and semantics.  A scoped derivation is a finite proof object built from named
rules and structural constructors.  Soundness is an induction principle on
that derivation.  Completeness is a separate metatheorem saying that the
chosen syntax presents all semantic homotopies.  This separation is exactly
what is lost when a global equality-normalization procedure is used as a
surrogate for a presentation.

\begin{theorem}[initial semantic factorization]
Let $\mathcal H$ be a topological groupoid over $X$, and let
\[
  \varphi:T\longrightarrow\mathcal H_1
\]
be continuous.  Assume that $\varphi$ preserves endpoints, reflexive traces,
reversal, and concatenation, and that every named scoped rewrite is sent to an
equality in $\mathcal H_1$.  Then there is a unique continuous map
\[
  \overline\varphi:G_{\mathcal P}\longrightarrow\mathcal H_1
\]
such that $\overline\varphi\,q=\varphi$.  It preserves all groupoid
operations, with composition interpreted on the final composable domain.
\end{theorem}

\begin{proof}
The rewrite assumptions say exactly that $\varphi$ is constant on the
scoped setoid.  The quotient universal property therefore defines
$\overline\varphi$ and gives its uniqueness.  Since $q$ is a quotient map and
$\varphi$ is continuous, the quotient criterion gives continuity.  The
operation-preservation equations hold on representatives by assumption and
descend through quotient extensionality.  On the final composable domain the
same argument uses the quotient of the explicit composable carrier, so no
product-quotient premise is introduced.
\end{proof}

The construction is compatible with the standard algebra of rewriting
\cite{ChurchRosser36,Newman42,KnuthBendix70,Squier87,GuiraudMalbos09}:
congruence, symmetry, transitivity, and critical-pair reasoning can be added
as named presentation rules.  Termination or confluence is not assumed by the
topological construction.  When they are used to construct a normalizer, the
normal-form completeness criterion above specifies the two additional semantic
checks that must be proved: geometric separation and scoped equality of equal
normal forms.  The finite-generator circle and torus reductions discharge
those checks explicitly.

The ambient semantics is deliberately modest.  The theorem uses ordinary
propositional equality, continuous interval paths, and endpoint-fixed
homotopies.  It does not claim a HoTT identity type, univalence, or a
proof-relevant equality groupoid.  This scope makes the comparison suitable
for a proof-irrelevant dependent type theory while retaining explicit
computational traces at the presented level \cite{MacLane98,Brown06}.

\section{Formalization boundary and artifact}

The mathematical statements above are independent of the implementation.  A
Lean 4.24.0 development with Mathlib 4.24.0 checks the construction in the
seven scoped-rewrite modules under
\LeanName{ComputationalPaths/Path/Topology/ScopedGeometricRewrite*.lean}.
The genuine torus certificate is in
\LeanName{ComputationalPaths/Path/Topology/TopologicalTorusScoped.lean}.
The root import hub is \LeanName{ComputationalPaths.lean}.  The principal
declaration layers are listed in Table~\ref{tab:declarations}.

\begin{table}[ht]
\centering
\small
\begin{tabular}{>{\raggedright\arraybackslash}p{0.29\linewidth}>{\raggedright\arraybackslash}p{0.63\linewidth}}
\toprule
Mathematical layer & Checked implementation declarations\\
\midrule
Presentation and soundness &
\LeanName{ScopedGeometricRewritePresentation},
\LeanName{ScopedRwEq}, \LeanName{ScopedRwEq.sound}\\
Quotient and final domain &
\LeanName{scopedEquivalent}, \LeanName{ScopedClass},
\LeanName{ScopedComposableClass}, \LeanName{scopedCompositionOnStrong}\\
Groupoid and topology &
\LeanName{ScopedStrongComposablePair},
\LeanName{scopedFinalTopologicalGroupoidCertificate}\\
Ordinary-pullback compatibility &
\LeanName{scopedProductCompatibility_iff_four_way},
\LeanName{scopedProductCompatibility_of_open_pair_map},
\LeanName{scopedProductCompatibility_of_compact_final_t2},
\LeanName{scopedProductCompatibility_of_discrete_arrow_and_final_domain}\\
Comparison and completeness &
\LeanName{toGeometricClass}, \LeanName{toGeometricStrongPair},
\LeanName{ComparisonFunctorCertificate},
\LeanName{comparisonHomeomorph_of_complete},
\LeanName{ScopedGeometricNormalFormCertificate},
\LeanName{geometricCompleteness_of_normalForm}\\
Universal completion &
\LeanName{universalPresentation},
\LeanName{universalScopedEquivalent_iff_totalEquivalent},
\LeanName{universalHomeomorph}\\
Fundamental-groupoid bridge &
\LeanName{RealizedFundamentalArrow},
\LeanName{RealizedFundamentalGroupoidCertificate},
\LeanName{universalRealizedFundamentalArrowHomeomorph}\\
Circle &
\LeanName{circleLoopRule}, \LeanName{circleTrace_normalizes},
\LeanName{circleLoopEquivInt},
\LeanName{circleBasedNormalFormCertificate},
\LeanName{circleFinalTopologicalGroupoidCertificate}\\
Actual topological torus &
\LeanName{TopologicalTorus.equivIntProd},
\LeanName{TopologicalTorus.standardLoop_homotopic},
\LeanName{TopologicalTorus.certificate}\\
\bottomrule
\end{tabular}
\caption{Theorem-to-declaration map for the artifact.}
\label{tab:declarations}
\end{table}

The realized fundamental-groupoid bridge identifies the quotient of realized
open paths with the range of the endpoint-and-homotopy code in Mathlib's
fundamental groupoid.  Its certificate includes explicit carrier membership
for identities, reversal, and composition.  The comparison certificate also
contains the continuous strong-pair map and final-composition equation, so the
formalized comparison is a groupoid morphism on the same final domain as the
mathematical theorem.

The root reproducibility command is \LeanName{lake build}; the focused module
commands are listed in the artifact README and can be run independently.

The checked modules contain no unfinished proofs and no custom axiom
declarations.  Declaration-level axiom inspection reports only Lean's standard
logical and quotient principles, namely propositional extensionality,
classical choice, and quotient soundness.  The Lean artifact is evidence for
the stated theorem package and complements the authors' current Lean preprint
\cite{deMoura2021,RamosEtAl2025LeanPreprint,RamosEtAl2025OmegaPreprint};
it is not used as a substitute for the
mathematical definitions or arguments in the main text.

\section{Conclusion}

Scoped computational paths provide a useful middle level between raw path
syntax and semantic fundamental-groupoid arrows.  The final-domain
construction is the key topological device: it gives a fully unconditional
groupoid multiplication while exposing, rather than hiding, the precise
product-quotient condition needed for the ordinary pullback topology.

The universal presentation shows that this is not merely a bookkeeping
construction: ordinary quotient-topologized fundamental groupoids arise as
the geometric completion of a scoped trace system.  The compact-Hausdorff and
discrete finite-presentation theorems give a positive ordinary-pullback class;
the effective normal-form criterion turns concrete reductions into geometric
completeness.  The finite-generator circle and genuine torus calculations then
show that the framework retains nontrivial topological information without
silently identifying scoped syntax with ambient equality.

\bibliographystyle{amsalpha}
\begingroup
\small
\bibliography{../refs}
\endgroup

\end{document}
